% Options for packages loaded elsewhere
% Options for packages loaded elsewhere
\PassOptionsToPackage{unicode}{hyperref}
\PassOptionsToPackage{hyphens}{url}
\PassOptionsToPackage{dvipsnames,svgnames,x11names}{xcolor}
%
\documentclass[
  french,
  letterpaper,
]{report}
\makeatother


\usepackage{longtable,booktabs,array}
\usepackage{calc} % for calculating minipage widths
% Correct order of tables after \paragraph or \subparagraph
\usepackage{etoolbox}
\makeatletter
\patchcmd\longtable{\par}{\if@noskipsec\mbox{}\fi\par}{}{}
\makeatother
% Allow footnotes in longtable head/foot
\IfFileExists{footnotehyper.sty}{\usepackage{footnotehyper}}{\usepackage{footnote}}
\makesavenoteenv{longtable}
\usepackage{graphicx}
\makeatletter
\newsavebox\pandoc@box
\newcommand*\pandocbounded[1]{% scales image to fit in text height/width
  \sbox\pandoc@box{#1}%
  \Gscale@div\@tempa{\textheight}{\dimexpr\ht\pandoc@box+\dp\pandoc@box\relax}%
  \Gscale@div\@tempb{\linewidth}{\wd\pandoc@box}%
  \ifdim\@tempb\p@<\@tempa\p@\let\@tempa\@tempb\fi% select the smaller of both
  \ifdim\@tempa\p@<\p@\scalebox{\@tempa}{\usebox\pandoc@box}%
  \else\usebox{\pandoc@box}%
  \fi%
}
% Set default figure placement to htbp
\def\fps@figure{htbp}
\makeatother


% definitions for citeproc citations
\NewDocumentCommand\citeproctext{}{}
\NewDocumentCommand\citeproc{mm}{%
  \begingroup\def\citeproctext{#2}\cite{#1}\endgroup}
\makeatletter
 % allow citations to break across lines
 \let\@cite@ofmt\@firstofone
 % avoid brackets around text for \cite:
 \def\@biblabel#1{}
 \def\@cite#1#2{{#1\if@tempswa , #2\fi}}
\makeatother
\newlength{\cslhangindent}
\setlength{\cslhangindent}{1.5em}
\newlength{\csllabelwidth}
\setlength{\csllabelwidth}{3em}
\newenvironment{CSLReferences}[2] % #1 hanging-indent, #2 entry-spacing
 {\begin{list}{}{%
  \setlength{\itemindent}{0pt}
  \setlength{\leftmargin}{0pt}
  \setlength{\parsep}{0pt}
  % turn on hanging indent if param 1 is 1
  \ifodd #1
   \setlength{\leftmargin}{\cslhangindent}
   \setlength{\itemindent}{-1\cslhangindent}
  \fi
  % set entry spacing
  \setlength{\itemsep}{#2\baselineskip}}}
 {\end{list}}
\usepackage{calc}
\newcommand{\CSLBlock}[1]{\hfill\break\parbox[t]{\linewidth}{\strut\ignorespaces#1\strut}}
\newcommand{\CSLLeftMargin}[1]{\parbox[t]{\csllabelwidth}{\strut#1\strut}}
\newcommand{\CSLRightInline}[1]{\parbox[t]{\linewidth - \csllabelwidth}{\strut#1\strut}}
\newcommand{\CSLIndent}[1]{\hspace{\cslhangindent}#1}

\begin{document}

\begin{table}

\caption{\label{tbl-menaces}Menaces perçues sur l'environnement et
responsables identifiés (\% de Oui)}

\centering{

\fontsize{12.0pt}{14.0pt}\selectfont
\begin{tabular*}{\linewidth}{@{\extracolsep{\fill}}lrr}
\toprule
 & {\bfseries {Alaotra}} & {\bfseries {Marovoay}} \\ 
\midrule\addlinespace[2.5pt]
\multicolumn{3}{l}{{\bfseries {Perception globale}}} \\[2.5pt] 
\midrule\addlinespace[2.5pt]
Considèrent que l'environnement est menacé & 82 % & 79 % \\ 
\midrule\addlinespace[2.5pt]
\multicolumn{3}{l}{{\bfseries {Principales menaces (parmi ceux qui déclarent menacé)}}} \\[2.5pt] 
\midrule\addlinespace[2.5pt]
Feux & 84 % & 95 % \\ 
Défrichage & 56 % & 52 % \\ 
Prélèvement (faune) & 40 % & 20 % \\ 
Prélèvement (flore) & 40 % & 21 % \\ 
Prospection minière & 29 % & 6 % \\ 
Nomadisme (passage non-résidents) & 16 % & 16 % \\ 
Zébu / pâturage & 20 % & 10 % \\ 
Tourisme & 2 % & 3 % \\ 
\midrule\addlinespace[2.5pt]
\multicolumn{3}{l}{{\bfseries {Principaux responsables}}} \\[2.5pt] 
\midrule\addlinespace[2.5pt]
L'ensemble des riverains & 87 % & 70 % \\ 
Les riverains migrants & 47 % & 37 % \\ 
L'État (Fanjakana) & 10 % & 8 % \\ 
Les braconniers & 30 % & 10 % \\ 
Le gestionnaire & 9 % & 3 % \\ 
Les nomades ou semi-nomades & 22 % & 17 % \\ 
Les ONG ou entreprises & 3 % & 4 % \\ 
Les chercheurs & 0 % & 2 % \\ 
Les touristes & 1 % & 2 % \\ 
\bottomrule
\end{tabular*}
\begin{minipage}{\linewidth}
\textbf{Source : Enquête auprès des OR 2025}\\
\end{minipage}

}

\end{table}%
\end{document}
